%% Reviewer response
\noindent\textbf{Response:}
We thank the reviewer for this helpful suggestion. We have added three
zero-centred summary plots for Tables~I, II, and IV. For each experimental
condition, we calculate
\[
\Delta=100\frac{C_{\mathrm{best\ baseline}}-C_{\mathrm{DVNDA}}}
{C_{\mathrm{best\ baseline}}}.
\]
Each plot uses a single forest-plot axis: the point shows the mean percentage
difference, the horizontal line shows its approximate 95\% confidence
interval, and the dashed vertical line denotes zero difference. Values to the
right of zero indicate a lower DVNDA cost; values to the left indicate a lower
cost for the best baseline. This compact, single-colour presentation avoids
visual clutter while showing both effect size and uncertainty. The Table~I
plot is placed in the main manuscript, whereas the Table~II and IV plots are
placed in the Supplementary Material. The complete tables are retained for
exact values and method-level comparisons.

%% Main-text insertion after Table I
\begin{figure}[t]
  \centering
  \includegraphics[width=\columnwidth]{figure/main_table_I_improvement.pdf}
  \caption{\textbf{Relative routing-cost difference between DVNDA and the best
  baseline on synthetic instances.} For every dynamic rate and problem size,
  $\Delta=100(C_{\mathrm{best\ baseline}}-C_{\mathrm{DVNDA}})/
  C_{\mathrm{best\ baseline}}$. Points to the right of zero indicate that
  DVNDA has the lower mean cost; points to the left indicate that the best
  non-DVNDA method in Table~\ref{tab:2} has the lower mean cost. Horizontal
  lines denote approximate 95\% confidence intervals
  obtained by an independent-sample delta method from the reported means and
  standard deviations ($n=100$ instances per condition). Row labels give the
  customer count and dynamic rate, separated by a vertical bar. An interval crossing
  zero does not establish a resolved difference under this approximation.}
  \label{fig:table1_improvement}
\end{figure}

%% Concise main-text interpretation
Figure~\ref{fig:table1_improvement} separates mean direction from uncertainty.
DVNDA has the lower mean cost in 9 of the 12 conditions, whereas the best
baseline has the lower mean cost in the three $\varphi=10\%$ conditions. At
$\varphi=10\%$, DVNDA is $1.96\%$--$2.77\%$ more costly than the best baseline.
From $\varphi=25\%$ onward, DVNDA is consistently less costly, with mean
reductions of $0.71\%$--$2.60\%$. However, all approximate 95\% intervals cross
zero. Thus, Table~I supports a consistent mean advantage at moderate and high
dynamic rates, but not a statistically resolved advantage from the available
summary statistics.

%% Supplementary Figure S1 (Table II)
\begin{figure}[t]
  \centering
  \includegraphics[width=\columnwidth]{figure/supp_table_II_improvement.pdf}
  \caption{\textbf{Relative routing-cost difference on Vienna road-network
  instances.} The signed percentage difference is
  $\Delta=100(C_{\mathrm{best\ baseline}}-C_{\mathrm{DVNDA}})/
  C_{\mathrm{best\ baseline}}$. Points to the right of zero indicate a lower
  DVNDA mean cost; points to the left indicate a lower mean cost for the best
  non-DVNDA method in Table~\ref{tab:zs}. Horizontal lines denote approximate
  95\% confidence intervals
  obtained by an independent-sample delta method from the reported means and
  standard deviations ($n=100$ instances per condition). Row labels give the
  customer count and dynamic rate, separated by a vertical bar.}
  \label{fig:supp_table2_improvement}
\end{figure}

On the Vienna instances, DVNDA has the lower mean cost in 9 of 12 conditions.
The three exceptions again occur at $\varphi=10\%$, where the best baseline is
$2.20\%$--$3.27\%$ less costly. At $\varphi=25\%$--$75\%$, DVNDA reduces the
mean cost by $0.93\%$--$5.87\%$; five of the corresponding nine intervals lie
entirely above zero. The real-road results therefore provide stronger evidence
of an advantage under moderate and high dynamism than the synthetic results.

%% Supplementary Figure S2 (Table IV)
\begin{figure}[t]
  \centering
  \includegraphics[width=\columnwidth]{figure/supp_table_IV_improvement.pdf}
  \caption{\textbf{Relative routing-cost difference under unseen spatial
  distributions.} Clustered and mixed instances contain 20 customers. The
  signed percentage difference is
  $\Delta=100(C_{\mathrm{best\ baseline}}-C_{\mathrm{DVNDA}})/
  C_{\mathrm{best\ baseline}}$. Points to the right of zero indicate a lower
  DVNDA mean cost; points to the left indicate a lower mean cost for the best
  non-DVNDA method in Table~\ref{fbfh}. Horizontal lines denote approximate
  95\% confidence intervals
  obtained by an independent-sample delta method from the reported means and
  standard deviations ($n=100$ instances per condition). Row labels give the
  spatial distribution and dynamic rate, separated by a vertical bar.}
  \label{fig:supp_table4_improvement}
\end{figure}

DVNDA has the lower mean cost in all eight cross-distribution conditions. The
reduction is $4.24\%$--$8.12\%$ on clustered instances and
$2.50\%$--$11.96\%$ on mixed instances. Five of eight intervals lie entirely
above zero, while the remaining three cross zero. These results support
cross-distribution robustness while preserving the uncertainty of individual
comparisons.

%% Statistical note for the authors
%% The manuscript tables contain marginal means and SDs but not paired
%% per-instance costs. The intervals above are therefore approximate. If the
%% paired raw costs are available, replace them with paired bootstrap 95% CIs
%% before final submission; the plotting code documents the replacement point.
